\section{Related Work}

Existing benchmark analysis approaches fall into three categories, none of which provide predictive coverage modeling.

\subsection{Benchmark Taxonomies}

Papers with Code~\cite{paperswithcode} and similar platforms organize benchmarks for browsing by task type (classification, generation, translation) and domain (vision, language, multimodal). While valuable for exploring existing resources, these taxonomies require manual review to determine hypothesis-specific suitability---researchers must read benchmark documentation to match evaluation frameworks to their validation needs. Our work automates this matching via design feature extraction and historical validation, predicting future suitability from benchmark characteristics rather than relying on manual review of organized listings.

\subsection{Citation-Based Analysis}

Semantic Scholar~\cite{semanticscholar} and Google Scholar enable benchmark popularity analysis through citation counts and co-citation networks. However, popularity does not equal coverage: ImageNet~\cite{imagenet} receives thousands of citations but is fundamentally unsuitable for sequence generation tasks due to metric incompatibility. These approaches measure usage frequency without extracting \textbf{why} benchmarks are suitable (design constraints). We extract design features (modality, metrics, task formulation) to predict benchmark-hypothesis compatibility independent of popularity.

\subsection{Benchmark Recommendation Systems}

Collaborative filtering and user-behavior similarity approaches recommend benchmarks based on historical usage patterns. These methods require prior usage data for each hypothesis type, making them unsuitable for novel hypothesis categories with no historical record. Our design-feature-based approach enables prediction without prior usage by leveraging inherent benchmark characteristics (data modality, evaluation metrics) that persist across publication cycles.

\subsection{Technical Foundations}

Our work builds on SciBERT~\cite{scibert} for scientific text understanding and SentenceBERT~\cite{sentencebert} for semantic similarity of benchmark descriptions. SciBERT's pre-training on scientific literature enables accurate citation context classification, while SentenceBERT embeddings capture semantic similarity enabling modality-driven patterns to emerge from text without requiring pre-defined categories. K-means clustering provides unsupervised discovery of coverage families, complemented by silhouette scoring for cluster quality validation.

\subsection{Positioning}

Unlike descriptive taxonomies (organize but don't predict), citation analysis (popularity $\neq$ coverage), or recommendation systems (require prior usage), we provide the first systematic method to predict future benchmark suitability from design features using historical train/test validation, reducing manual selection from weeks to minutes.
